\section{When geometric movement composes with computational work}
\label{work:section}

The resource inequalities in Section~\ref{res:section} become work lower
bounds only under an explicit cost premise. The point requiring proof is
that the rank paying for curvature is also the rank paying for movement.
A larger ambient or restricted barrier parameter cannot be substituted
for that rank in a movement lower bound.

\subsection{The same aggregate certificate in both estimates}
Consider the full rows of $C=\prod_{a=1}^b\B_2^{s_a}$, with
$s_a\geq2$, and globally bi-$C^1$ selections as in
\eqref{eq:full-rows}. Suppose each dual row is a genuine certificate on
the whole feasible affine slice. Choose positive weights $\lambda_a$ and
a simultaneous contact $u$. The objective and its certificate are
\[
 \ell(x)=\sum_a\lambda_a u_a^Tx_a,\qquad
 \ell_*=\sum_a\lambda_a,\qquad
 S_i=\sum_a\lambda_aY_i^a(u_a).
\]
Write $n=\sum_a(s_a-1)$, $p_i=\jrank X_i(u)$,
$q_i=\jrank S_i$ and $Q=\sum_iq_i$. For simple factors with rank $r_i$
and Peirce constant $a_i$, the mixed-channel bound gives
\begin{equation}\label{work:same-rank}
 n\leq\sum_{i\text{ non-ray}}a_i p_iq_i,
 \qquad p_i+q_i\leq r_i.
\end{equation}
To see this directly, the weighted slack has mixed derivative
$-\sum_a\lambda_a\ip{h_a}{k_a}$, a nondegenerate form on an
$n$-dimensional tangent space. At contact, positivity makes $X_i(u)$
and $S_i$ complementary. The differentiated mixed pairing factors through
the cross Peirce space between their supports, of dimension $a_ip_iq_i$.
The ranks add to at most $r_i$ by complementarity; rank subadditivity
proves~\eqref{work:same-rank}. In real PSD notation this is the space
$\operatorname{Hom}(\operatorname{Ran}S_i,\operatorname{Ran}X_i)$;
its transpose is not a second independent channel.
The whole-slice certificate identity allows
Theorem~\ref{thm:movement-minor} to use exactly this $Q$.

There is also an explicit factor-aware consequence for a common Jordan
family with Peirce constant $a>0$, orders at most an integer $R\geq2$,
and at most an integer $L\geq1$ non-ray factors. With $\mathcal C_R$ as in~\eqref{res:nullity-envelope},
set
\begin{equation}\label{work:qmin}
 Q_{\min}(n;L,a,R)=\min\{z\in\mathbb Z_{\geq0}:
                              a\mathcal C_R(L,z)\geq n\}.
\end{equation}
It exists only if $n\leq aL\lfloor R^2/4\rfloor$; otherwise there is
no such factorization. In the feasible case every aggregate certificate
above has $Q\geq Q_{\min}$. Optional rays can contribute to $Q$ but
not to capacity, and do not affect the assertion. For the non-ray rank
$Q_0$, Cauchy--Schwarz gives
$n\leq a(RQ_0-Q_0^2/L)$, whence also
\begin{equation}\label{work:continuous-rank}
 Q\geq\frac L2\left(R-\sqrt{R^2-\frac{4n}{aL}}\right).
\end{equation}
Thus at high curvature loading the movement coefficient is of order
$\sqrt{LR}$, even if the number of factors is only $L$.
For an exact fixed non-ray rank $Q_0\leq LR$, its untruncated maximum
capacity is
$a[(L-t)u(R-u)+t(u+1)(R-u-1)]$ with $Q_0=uL+t$.
For a rank \emph{budget}, the truncation in $\mathcal C_R$ is necessary:
capacity is maximized by stopping at $L\lfloor R/2\rfloor$.
All these envelopes are integer consequences of curvature, not claims
that each balanced rank vector integrates to a lift.

These bounds can also be imposed at every stage of the sequential
compression in Theorem~\ref{thm:smooth-product-rank}. For a dictionary
$\mathscr K$ whose non-ray faces retain the common Peirce constant $a$
and order bound $R$, the same proof selects a simultaneous contact with
\begin{equation}\label{work:sequential-budget}
 Q\geq\sum_{a'=1}^b
 \max\{\kappa_{\mathscr K}(s_{a'}),
           Q_{\min}(s_{a'}-1;L,a,R)\}.
\end{equation}
There are at most $L$ non-ray faces at every stage, and each compressed
full row has mixed rank $s_{a'}-1$; newly exposed ranks add by
Lemma~\ref{lem:join}. These facts prove the displayed bound. It applies
to each fixed classical Hermitian family and to a fixed spin family;
Albert faces also retain Peirce constant eight when non-ray. If these factors also have real dimension at most $d$, the proved bound
$h_d(s_{a'})$ of Section~\ref{sec:orbits} is an alternative one-row
bound and can be combined with $Q_{\min}$ in the same way. For mixed families use the adaptive
face ranks and constants in that section, instead of assigning them an
artificial common constant.

\subsection{A conditional total-work theorem}
Specialize to real PSD blocks $2\leq r_i\leq R$ and the standard
restricted barrier $F=-\sum_i\log\det X_i$. A method starts at a fixed
interior reference $x^0$. Every counted round moves intrinsic distance
at most $B>0$, and every round is charged at least $a_0J_\omega$ work,
where $a_0>0$ and $J_\omega=\sum_i r_i^\omega$. The charge is an
assumption about what the method must process or output in each round.
It is not implied by access to sparse entries or quantum oracles.

\begin{theorem}[Work under a fresh per-round charge]\label{work:total}
For the support objective and certificate above, let $\Delta>0$ be
\eqref{eq:movement-delta} with all weights $\alpha_i=1$. If
$0<\varepsilon<\Delta$ and $0\leq\omega\leq5/2$, every such method
reaching objective gap at most $\varepsilon$ has total charged work
\begin{equation}\label{work:total-bound}
 \mathcal W\geq\frac{a_0}{B}\frac{25}{6}\sqrt{\frac53}
 \frac{n^{3/2}}{R^{(5-2\omega)/2}}
 \log\frac{\Delta}{\varepsilon}.
\end{equation}
If each round instead costs at least $a_0M$, where
$M=\sum_i d_{r_i}$, then
\begin{equation}\label{work:coordinate-total}
 \mathcal W\geq\frac{a_0}{B}\frac{25}{12}\sqrt{\frac53}
 (1+1/R)\frac{n^{3/2}}{\sqrt R}\log\frac{\Delta}{\varepsilon}.
\end{equation}
A round of at most an integer $m\geq1$ feasible chords of starting local norm at
most $0<\eta<1$ satisfies the distance premise with
$B=m\log(1/(1-\eta))$.
\end{theorem}
\begin{proof}
Apply the algebraic proof of Theorem~\ref{res:holder} to
$p_i\leq r_i-q_i$ and~\eqref{work:same-rank}, taking $\theta=1/2$.
It gives
\[
 J_\omega\sqrt Q\geq\frac{25}{6}\sqrt{\frac53}
                    \frac{n^{3/2}}{R^{(5-2\omega)/2}}.
\]
Theorem~\ref{thm:movement-minor} bounds the endpoint distance below by
$\sqrt Q\log(\Delta/\varepsilon)$. The triangle inequality therefore
requires at least this distance divided by $B$ rounds. Multiplication
by the stipulated charge proves~\eqref{work:total-bound}.
The inequality $M\geq(1+1/R)J_2/2$ proves the coordinate version.
The chord estimate is the classical self-concordant Hessian comparison
already used in Corollary~\ref{cor:movement-rounds}.
\end{proof}

\subsection{Attainment in specified output models}
For homogeneous source dimensions $s$ suppose the integers allow
$p=3R/5$, $c=2R/5$, $h=s/p$ and $L=bh/c$.
The packed lift has $H=bh$ columns and full order-$R$ blocks.
For unit source directions $v_a$, maximize $\sum_av_a^Tx_a$.
The full fiber-centered path of Proposition~\ref{prop:cm-grouped} has
\[
 x_a=r(\tau)v_a,\quad
 r(\tau)=\frac{\tau}{\sqrt{h^2+\tau^2}+h},\qquad
 D_j=q(\tau)I_c,
 \quad q(\tau)=\frac{2}{\sqrt{h^2+\tau^2}+h}.
\]
At its support contact, the $c$ dual column vectors in each block are
linearly independent because their allocation coordinates are different
basis vectors. Hence $q_i=c$ and $Q=H$. Its resource product is exactly
\begin{equation}\label{work:matching-product}
 J_\omega\sqrt H=\frac{25}{6}\sqrt{\frac53}
                 \frac{(bs)^{3/2}}{R^{(5-2\omega)/2}}.
\end{equation}
Its distance from the analytic center to accuracy $\varepsilon$ is at
least $\sqrt H[\log(b/(2\varepsilon))]_+$, and its central arc has
length $O(\sqrt H\log(b/\varepsilon))$ for
$0<\varepsilon\leq b/4$. Thus bounded-distance checkpoints attain the
movement order with the resource ratio $(s/(s-1))^{3/2}$ to
\eqref{work:total-bound}. The instance-dependent reference scale in
that theorem is retained; no uniform scale over all lifts is asserted.

For a concrete dense-output upper cost, preprocess the fixed objective
columns $V_j$ and their Gram matrices $G_j=V_j^TV_j$ once. At each
checkpoint the entire lifted point is
\[
 W_j=rV_j,\qquad S_j=r^2G_j+qI_c.
\]
Writing all ambient cone entries therefore takes $O(M)$ exact real
arithmetic operations and writes per checkpoint, including the fixed
identity blocks if the output contract requires them. This proves a
matching \emph{postprocessing and serialization} bound for that path,
with preprocessing and scalar evaluation charged separately. It does
not require a fresh Gram computation at every round. An explicit
bounded-distance schedule also avoids a hidden optimization subproblem:
from $\tau=0$ to $h$, use increments $O(h/\sqrt H)$; subsequently use
constant increments $O(1/\sqrt H)$ in $\log\tau$ until
$\tau=H/\varepsilon$. The squared log-parameter speed is
$H(1-h/\sqrt{h^2+\tau^2})\leq H$, and before $h$ the ordinary speed
is at most $\sqrt{H/2}/h$. Also $b(1-r(\tau))\leq H/\tau$.
This verifies both its step budget and its endpoint accuracy. Square-root
and exponential evaluations, or their chosen precision implementations,
are separate scalar costs; there is no hidden dense linear solve here.

A charge of one record per factor ($J_0$), or $r_i$ records per factor
($J_1$), is a different output model. If those records consist of the
public scalars and references to preprocessed fixed directions and Gram
matrices, they can be refreshed in $O(J_0)$ or $O(J_1)$ word operations.
Equation~\eqref{work:matching-product} then describes matching charged
record work for this representation. It does not claim that all $bs$
projected coordinates, or all dense factors, can be produced in that
cost. Conversely, the full-coordinate model really does require fresh
writes of every output entry in every counted round. Counting only free
coordinates instead leads to~\eqref{res:free-optimum}, whose optimizer
is grouped Schur rather than $3:2$ packing.

Finally, none of these fresh charges follows from one-shot query
complexity. Fixed input data can be preprocessed and reused, and public
constant cone blocks can be added without changing the problem. For
example, let $v_{aj}=\sigma_{aj}/\sqrt s$ with hidden signs
$\sigma_{aj}\in\{-1,1\}$. One phase-oracle query prepares
\[
 \frac1{\sqrt{bs}}\sum_{a,j}\sigma_{aj}|a,j\rangle.
\]
Every nonzero projected center and projected predictor on the displayed
path is a public scalar multiple of this same vector. Fresh preparation
costs one hidden query per checkpoint, independently of the number of
PSD factors; it does not cost $\Omega(L)$ queries. This says nothing
about auxiliary Gram-state preparation, final classical readout, or other
solver subroutines. It demonstrates why a one-shot lower bound cannot be
multiplied by a movement count without a further temporal or output
premise.
